%
% File: abstract.tex
% Author: V?ctor Bre?a-Medina
% Description: Contains the text for thesis abstract
%
% UoB guidelines:
%
% Each copy must include an abstract or summary of the dissertation in not
% more than 300 words, on one side of A4, which should be single-spaced in a
% font size in the range 10 to 12. If the dissertation is in a language other
% than English, an abstract in that language and an abstract in English must
% be included.

\chapter*{R\'{e}sum\'{e}}
\initial{D}e nos jours, dans les entreprises, nous sommes amen\'{e}s \`{a} g\'{e}rer plusieurs types de documents (contrats, factures, bulletins de paie, CV). L'extraction de donn\'{e}es issues de cette vari\'{e}t\'{e} de documents repr\'{e}sente un d\'{e}fi majeur pour les syst\`{e}mes d'information. En effet, les solutions actuelles manquent d'efficacit\'{e}~: les approches bas\'{e}es sur des r\`{e}gles strictes sont rigides, tandis que les mod\`{e}les d'intelligence artificielle sacrifient souvent la pr\'{e}cision sur les valeurs num\'{e}riques et le contexte positionnel. Pour pallier \`{a} cela, notre recherche met en \oe{}uvre une m\'{e}thodologie hybride innovante qui combine l'extraction par templates, la reconnaissance d'entit\'{e}s nomm\'{e}es (NER) et les techniques de reconnaissance de motifs. Dans un cadre agile, notre syst\`{e}me int\`{e}gre un m\'{e}canisme d'apprentissage continu qui am\'{e}liore significativement la pr\'{e}cision d'extraction au fil des it\'{e}rations m\'{e}tier, en s'adaptant aux nouveaux formats de documents sans n\'{e}cessiter de r\'{e}entra\^{\i}nement lourd. Notre prototype, int\'{e}gr\'{e} comme module \`{a} DOLibaar, a r\'{e}duit la saisie manuelle de 80~\% et atteint plus de 95~\% de pr\'{e}cision d'extraction sur un p\'{e}rim\`{e}tre de 10\,000 documents tests. En comparaison, une approche par expressions r\'{e}guli\`{e}res seules plafonne \`{a} 72~\% de pr\'{e}cision, tandis qu'un mod\`{e}le LLM g\'{e}n\'{e}rique non fine-tun\'{e} atteint \`{a} peine 81~\% sur le m\^{e}me jeu de donn\'{e}es, d\'{e}montrant ainsi la sup\'{e}riorit\'{e} de notre m\'{e}thode hybride. Il s'int\`{e}gre dans les syst\`{e}mes m\'{e}tier via des formats standards (JSON, XML), pour un temps de traitement moyen inf\'{e}rieur \`{a} 500~ms par fichier. Bien que la gestion des documents scann\'{e}s de faible qualit\'{e} reste une limite identifi\'{e}e, notre architecture modulaire ouvre la voie \`{a} l'ajout d'un pr\'{e}traitement OCR, consolidant ainsi l'industrialisation de la solution.

\paragraph{Mot cl\'{e}s : Extraction automatique de donn\'{e}es, Documents administratifs, Reconnaissance d'entit\'{e}s nomm\'{e}es (NER), M\'{e}thodologie hybride, Apprentissage continu, DOLibaar.}


\clearpage

